% \section{Appendix}
% \label{sec:append}


\section{Shifting step}\label{sec:shifting}
%\begin{restatable}{lem}{shiftinglemmaapx}
\begin{lem}
	\label{lem:shifting}
	Let a given packing for a knapsack packing problem with input set \(\C\) of items be $\mathcal{P}$ and \(r\) be a total order on these objects and \(w_i\) be the profit associated with the \(i\)-th object. Let $\mathcal{C}_{\mathcal{P}}$ be the set of the items selected by $\mathcal{P}$ and we have $w(\mathcal{P}) = \sum_{C_{i}\in \mathcal{C}_{\mathcal{P}}} w_{i}$. Define $\rhoi{k} = h(\eps, k)$ where $h: \mathbb{R} \rightarrow \mathbb{R}$ is a given decreasing function. For each integer \(k\geqslant 1\), define the following sets of items:
	\(
	\mathcal{C}_k = \{C_i \in \mathcal{C} \mid r_{i} \in (\rho_{k},\rho_{k-1}]\}.
	\)
	Then there exists a \(\tau \leqslant \frac{1}{\eps}\) such that \(w(\C_{\tau}\cap\C_{\mathcal{P}})\leqslant\eps w(\mathcal{P})\).
\end{lem}
%\end{restatable}
\begin{proof}
We will partition the input set of $\mathcal{C}$ into large, medium, and small items using a partition function $\rho_k$ which we will design retrospectively after considering all the properties it needs to satisfy. The methodology is as follows, we partition the input $\mathcal{C}$ into several sets \(\mathcal{C}_k\)'s, and then select an appropriate index \(j \in [1/\eps]\) such that the sum of weights of all items from $\mathcal{P}$ in \( \mathcal{C}_j \leqslant \eps\cdot w(\mathcal{P})\) for a given constant $\eps>0$. We will declare this to be the medium items set $\mathcal{C}_M$ and discard it from $\mathcal{C}$.

Let $\eps>0$ be a positive real constant. We can assume that $\mathcal{P}$ has more than $\frac{1}{\eps}$ items. Otherwise, it is trivial.
%we can enumerate all the sets with at most \(\frac{1}{\eps}\) items. Among the selected items, we note \(r\) \ of the smallest item and let it be denoted as \(r_{\text{min}}\). For this selected set of items, we partition them with all the items categorized as large items and both the set of medium and small items as empty. Formally we define \(\mathcal{C}_L = \{C_{i} \in \mathcal{C} \mid r_{i}\geq r_{\text{min}}\}\) which contains all the items and then set \(\mathcal{C}_M\) and \(\mathcal{C}_S\) are empty.
Now, we define $\rho_0=1$, and $\rhoi{k} := h(\eps, k) \cdot \rho_{k-1}$ where $h: \mathbb{R} \rightarrow \mathbb{R}$ is the given decreasing function. For each integer $k \geqslant 1$ we define the following sets of items:
\begin{align*}
\mathcal{C}_k = \{C_i \in \mathcal{C} \mid r_{i} \in (\rho_{k},\rho_{k-1}]\}.
\end{align*}

Let $\tau$ be the smallest index such that the total weight, $w(\mathcal{C}_{\tau})$, of the items in $S_{\tau}$ is at most $\eps w(\mathcal{P})$. Observe that $\tau \leqslant \frac{1}{\eps}$ due to the pigeonhole principle. Partition the set $\mathcal{C}$ of items into three groups: the $\textit{large}$ items $\mathcal{C}_L = \{C_{i} \in \mathcal{C} \mid r_{i}>\rho_{\tau-1}\}$, the $\textit{medium}$ items $\mathcal{C}_M = \{C_{i} \in \mathcal{C} \mid \rho_{\tau} < r_{i} \leqslant \rho_{\tau-1}\}$ and the $\textit{small}$ items $\mathcal{C}_S = \{C_{i} \in \mathcal{C} \mid r_{i} \leqslant \rho_{\tau}\}$. Similarly, define sets $\mathcal{C}_L^* = \mathcal{C}_L \cap \mathcal{P}$, $\mathcal{C}_M^* = \mathcal{C}_M \cap \mathcal{P}$, and $\mathcal{C}_S^* = \mathcal{C}_S \cap \mathcal{P}$. Note that $\mathcal{C}_M^* = \mathcal{C}_{\tau}$, so the total weight of the medium items is at most $\eps\cdot w(\mathcal{P})$.
\end{proof}


\section{\textsf{PTAS} for $d$-dimensional hyperspheres}
In this section we prove \Cref{thm:d-dim-ptas}.
	In order to extend our techniques to \(d\)-dimension, we start with first creating partitions over the \(d\)-dimensional hyperspheres, similar to Lemma \ref{lem:create-gap},  by setting \(\eps^{j}=(\eps^{j-1})^{12d}\) and assume that medium hyperspheres contribute negligible profit to $\OPT$.
	
	Further, we notice that each large hypersphere has a volume of at least \( \dfrac{\pi^{\frac{d}{2}}}{\Gamma\left(\frac{d}{2} + 1\right)} \cdot \epsl^d
	\). Thus, there can be at most  \(\left(\frac{\pi^{\frac{d}{2}}}{\Gamma\left(\frac{d}{2} + 1\right)} \cdot \epsl^d\right)^{-1}\)  
 %\(\dfrac{1}{\frac{\pi^{\frac{d}{2}}}{\Gamma\left(\frac{d}{2} + 1\right)} \cdot \epsl^d}\) 
 large hyperspheres in any feasible solution. Using the techniques as in \Cref{subsec:Large_Guess}, we can guess and verify the coordinates of the large hyperspheres in \(\mathcal{C}_{L}^*\) in polynomial time. Also observe that the guessed values $\tilde{x}_{i}^{(1)},\tilde{x}_{i}^{(2)},\dots,\tilde{x}_{i}^{(d)}$
	yield an estimate for $\hat{x}_{i}^{(1)},\hat{x}_{i}^{(2)},\dots,\hat{x}_{i}^{(d)}$ up to a (polynomially small) error of $\eps/n$.
	
	In order to find a feasible packing of the small hyperspheres we will replicate techniques from \Cref{subsec:placing_small}.  Define \(\epsc:=\epsl^{6d}\)(i.e, \(\epss=\epsl^{12d}=\epsc^2\)). To subdivide the knapsack into \(1/\epsc^d\) $d$-dimensional grid cells of length \(\epsc\). We denote a grid cell by \(G_{\ell^{(1)},\dots,\ell^{(d)}}:=[\ell^{(1)}\cdot\epsc,(\ell^{(1)}+1)\cdot\epsc)\times\dots\times[\ell^{(d)}\cdot\epsc,(\ell^{(d)}+1)\cdot\epsc)\). We define the set of grid cells by \(\G:=\{G_{\ell^{(1)},\dots,\ell^{(d)}}:\ell^{(1)},\dots,\ell^{(d)}\in\{0,1,\dots,1/\epsc-1\}\}\). To proceed further, we first prove structural properties for the packing of the hyperspheres in the grids. For any set \(\mathcal{S}\) of circles or grid cells, we define \(\area(\mathcal{S})\) to be the total volume of the elements in \(\mathcal{S}\).
	\begin{lem}\label{lem:small-hypersphere-grid-cells}
		Let \(\mathcal{C}_{cut}\subseteq\C_{S}^{*}\) be the set of all small hyperspheres in \(\C_{S}^{*}\) that intersect more than one grid cell. We have that \(\area(\mathcal{C}_{cut})\leqslant4d\epss/\epsc\leqslant d\eps\epsl^d/2^{5d-3}\)
	\end{lem}
	\begin{proof}
		For each of the \(i\)-th dimension for \(1\leqslant i \leqslant d\), and for each of the gridline \(\ell\) in this dimension, each \(C\in\C_{cut}\) must lie in a strip with the \(i\)-th dimension being \(4\epss\) and all other dimensions being 1. Totally, there are \(d/\epsc\) gridlines. Thus, \(\area(\C_{cut})\leqslant4d\epss/\epsc=4d\epsl^{6d}\leqslant d\eps\epsl^d/2^{5d-3}\). The last inequality follows from the fact \(\epsl\leqslant \eps\leqslant 1/2\).
	\end{proof}
	We partition \(\G\) into \emph{white}, \emph{black}, and \emph{gray} cells with respect to the position of the large hyperspheres similarly as the case of circles. We give bounds similar to \Cref{lem:gray-circles} and \Cref{lem:white-circles} in the case of $d$-dimensions for \emph{gray} and \emph{white} cells.
	\begin{lem}\label{lem:gray-circles-d}
		The total volume of gray cells \(\area(\G_{g})\) is at most \(\dfrac{(d+1)(3^d-1)\eps\epsl^d}{2^{5d-3}}\). We can compute \(\G_g\) in polynomial time.
	\end{lem}
	\begin{proof}
		Consider a large hypersphere \(C_L\) with radius \(r_L\geqslant\epsl\). Let \(\G_{g, C_L},\G_{b, C_L}\) be the set of gray and black grid cells that overlap with \(C_L\). Note that, for \(C_L\) we only have an estimate \(\tilde{x}_{C}^{(1)},\dots,\tilde{x}_{C}^{(d)}\) of the coordinate of its center up to a \(\emph{polynomially}\) small error of \(\eps/n\).
		
		First we claim that all the cells intersected by a hypersphere centered at  \((\tilde{x}_{C}^{(1)},\dots,\tilde{x}_{C}^{(d)})\) of radius $r_L- (d+1)\epsc$ are in $\G_{b, C_L}$. 
		For contradiction, assume there is a gray cell $S_g$ intersected by a circle $C$ centered at  \((\tilde{x}_{C}^{(1)},\dots,\tilde{x}_{C}^{(d)})\) of radius $r_L- (d+1)\epsc$.
		Then, the distance of any point in $S_g$ from the center of the circle is at most $r_L- (d+1)\epsc +\sqrt{d}\epsc < r_L-\eps/n$. Thus, $S_g$  is completely contained within $C$.
		
		Similarly, we can show that the cells that are outside of a circle centered at  \((\tilde{x}_{C}^{(1)},\dots,\tilde{x}_{C}^{(d)})\) of radius $r_L+ (d+1)\epsc$ are not in $\G_{b, C_L}$.
		
		Hence, 
		\begin{align*}
		\area(\G_{b, C_L}) &\geqslant \frac{\pi^{\frac{d}{2}}}{\Gamma\left(\frac{d}{2} + 1\right)} \cdot (r_L-(d+1)\epsc)^d \\
		\area(\G_{g, C_L}) &\leqslant \frac{\pi^{\frac{d}{2}}}{\Gamma\left(\frac{d}{2} + 1\right)} \cdot (r_L+(d+1)\epsc)^d - \frac{\pi^{\frac{d}{2}}}{\Gamma\left(\frac{d}{2} + 1\right)} \cdot (r_L-(d+1)\epsc)^d \\ 
		&\leqslant \frac{\pi^{\frac{d}{2}}}{\Gamma\left(\frac{d}{2} + 1\right)} \cdot 2(d+1)((3/2)^d-(1/2)^d)\epsc r_{L}^{d-1}
		\end{align*}
		\begin{comment}
		
		$\area(\G_{b, C_L}) \geqslant \dfrac{\pi^{\frac{d}{2}}}{\Gamma\left(\frac{d}{2} + 1\right)} \cdot (r_L-(d+1)\epsc)^d$, and\\ $\area(\G_{g, C_L}) \leqslant \dfrac{\pi^{\frac{d}{2}}}{\Gamma\left(\frac{d}{2} + 1\right)} \cdot (r_L+(d+1)\epsc)^d - \dfrac{\pi^{\frac{d}{2}}}{\Gamma\left(\frac{d}{2} + 1\right)} \cdot (r_L-(d+1)\epsc)^d\leq \dfrac{\pi^{\frac{d}{2}}}{\Gamma\left(\frac{d}{2} + 1\right)} \cdot 2(d+1)((3/2)^d-(1/2)^d)\epsc r_{L}^{d-1}$. 
		\end{comment}
		The inequality follows from the fact, \(r_{L}/2\geqslant(d+1)\epsc\)\\
		
		Hence we can bound the ratio of \(\frac{\area(\G_{g, C_L})}{\area(\G_{b, C_L})}\) as
		\[
		\begin{aligned}
		\frac{\area(\G_{g, C_L})}{\area(\G_{b, C_L})} &\leqslant \frac{\frac{\pi^{\frac{d}{2}}}{\Gamma\left(\frac{d}{2} + 1\right)} \cdot 2(d+1)((3/2)^d-(1/2)^d)\epsc r_{L}^{d-1}}{\frac{\pi^{\frac{d}{2}}}{\Gamma\left(\frac{d}{2} + 1\right)} \cdot (r_L-(d+1)\epsc)^d} \\
		&\leqslant \frac{2(d+1)((3/2)^d-(1/2)^d)\epsc r_{L}^{d-1}}{(r_L/2)^2} \\
		&\leqslant \frac{2(d+1)(3^d-1)\epsc}{r_L} \leqslant \frac{2(d+1)(3^d-1)\epsc}{\epsl} \\
		&= 2(d+1)(3^d-1)\epsl^{6d-1} \leqslant \frac{(d+1)(3^d-1)\eps\epsl^d}{2^{5d-3}}
		\end{aligned}
		\]
		
		Inequalities follow from the fact that $\epsl \leqslant \eps \leqslant 1/2$ and $r_L \geqslant \epsl$.  The last inequality uses the fact that $\eps \leqslant 1/2$.
		Therefore, $\sum_{C \in \C_{L}^{*}} {\area(\G_{g, C_L})} \leqslant \bigg(\dfrac{(d+1)(3^d-1)\eps\epsl^d}{2^{5d-3}}\bigg) \cdot \sum_{C \in \C_{L}^{*}}  {\area(\G_{b, C_L})} \leqslant \dfrac{(d+1)(3^d-1)\eps\epsl^d}{2^{5d-3}}$.
	\end{proof}
	\begin{lem}\label{lem:white-circles-d}
		The total volume of the white grid cells \(\area(\G_w)\) is at least \(\frac{\epsl^d}{2^d}\), We can compute \(\G_w\) in polynomial time.
	\end{lem}
	\begin{proof}
		Similarly arguing as in \Cref{lem:white-circles}, we can show that no hypersphere can intersect a $d$-dimensional bin of side length \(\epsl/4\) at each corner. Hence the total volume of white grid cells must be at least \(2^d\cdot(\epsl/4)^d=(\epsl/2)^d\)
	\end{proof}
	
	We note here that the total volume of gray cell and \(\C_{cut}\) is at most \(d\eps\epsl^d/2^{5d-3}+\dfrac{(d+1)(3^d-1)\eps\epsl^d}{2^{5d-3}} \leqslant \eps (\epsl/2)^d\). Hence similar to \Cref{lem:small-circles-white-grid-cells} we state the following lemma for $d$-dimensional case also using \Cref{lem:small-hypersphere-grid-cells}, \Cref{lem:gray-circles-d}, and \Cref{lem:white-circles-d}.
	\begin{lem}
		\label{lem:small-circles-white-grid-cells-d}
		There
		is a set $\C'_S \subseteq \C^*_S$ of small hyperspheres with
		$p(\C'_S)\geqslant (1-\eps) p(\C^*_S)$
		such that there is a packing for $\C'_S$ using the grid cells in $\G_w$ only.
	\end{lem}
	We extend \Cref{lem:grid-resource-augmentation} for the $d$-dimensional case.
	\begin{lem} \label{lem:grid-resource-augmentation-d}
		There is a set of small hyperspheres $\C''_{S}\subseteq\C'_{S}$ such that $p(\C''_{S})\geqslant(1-\eps)p(\C'_{S})$
		and it is possible to place the hyperspheres in $\C''_{S}$ non-overlappingly
		inside $|\G_{w}|$ $d$-dimensional knapsacks of size $(1-\eps)\epsc \times\cdots \times(1-\eps)\epsc$
		each. \end{lem}
	\begin{proof}
			Consider a $d$-dimensional knapsack in $\G_{w}$ of side length $\epsc$.
		Similar to Lemma \ref{lem:grid-resource-augmentation-1}, by removing all circles intersected by a random strip of width $\eps \cdot \epsc$ in all the $d$ directions, we incur a loss of at most $d(d r/\epsc + \eps)$ fraction of optimal profit. As $r\le \epss$, we incur a loss of at most $O(\eps)$ fraction of profit as $\epss = \epsc^2 = \epsl^{12d}$.
	\end{proof}
	We further extend \Cref{lem:structured-packing} to $d$-dimensional case using \Cref{lem:grid-resource-augmentation-d}
	\begin{lemma}\label{lem:structured-packing-d}
		In polynomial time, we can compute a set of grid cells $\G_w$ such that no grid cell in $\G_w$ intersects with any
		hypersphere $C_i \in\C_{L}^{*}$ for any legal placement of $C_i$. Moreover, there is a set of small circles $\C_S \subseteq \C_{S}^{*}$ such that $w(\C_{S}) \geqslant (1-\eps)w(\C_{S}^{*})$ and
		the circles in $\C_S$ can be packed non-overlappingly inside $|\G_w|$ grid cells of size $(1-\eps) \epsc \times\dots\times(1-\eps) \epsc$ each.
	\end{lemma}
	%Using it, we compute an approximation to the set $\C_S$ via the following result from \cite{MiyazawaWAOA}.
		Using it, we compute an approximation to the set $\C_S$ via an extension of our result from \Cref{thm:fat-ra} to $d$-dimensions and setting fatness $f$ to 1 for spheres.
	
	
	We pack the computed hyperspheres into our grid cells $\G_w$, denote them by $\C'_S$. In particular, they do not intersect any of the
	large hyperspheres in $\C_{L}^{*}$ in any legal placement of them. Our solution (corresponding to the considered guesses) consists of $\C_{L}^{*}\cup \C'_S$.

In order to further improve the precision of the large hyperspheres, we replicate techniques from \Cref{subsec:Improve_Large} which completes the proof of \Cref{thm:d-dim-ptas}.
